\documentclass[11pt]{article}
\usepackage[margin=1in]{geometry}
\usepackage{amsmath,amssymb,amsthm}
\usepackage{hyperref}

\newtheorem{theorem}{Theorem}
\newtheorem{proposition}{Proposition}
\newtheorem{lemma}{Lemma}
\newtheorem{conjecture}{Conjecture}
\newtheorem{remark}{Remark}
\newcommand{\mc}{\operatorname{MaxCut}}
\newcommand{\dd}{d}
\newcommand{\E}{\mathbb{E}}

\title{Exchange Certificates for an Inductive Approach to Triangle-Free Graph Bipartization}
\author{Jinji Li}
\date{Version 1.0 draft -- September 2026}

\begin{document}
\maketitle

\begin{abstract}
For a graph $G$, let
\[
\dd(G)=|E(G)|-\mc(G)
\]
be the minimum number of edges whose deletion makes $G$ bipartite.
A classical conjecture of Erd\H{o}s asserts that every triangle-free graph
on $N$ vertices satisfies $\dd(G)\le N^2/25$; for $N=5k$ this becomes
$\dd(G)\le k^2$.

We study an induction strategy based on deleting five vertices at a time,
with target
\[
\exists X\subseteq V(G),\ |X|=5:\qquad
\dd(G)-\dd(G-X)\le 2k-1.
\]
The paper develops three complementary pieces of evidence for this route:
exact balanced and mixed extension certificates, a coordinated-flip
inequality that converts local recolorings into a selection criterion, and
an exact 3-for-3 exchange audit in the difficult $n=15$ regime.

A stronger fixed-$X$ gain statement suggested by early experiments is
false on explicit realizable residual configurations.  The failure does
not contradict the induction target: exact audits show that changing the
five-set repairs the obstruction.  In the saved $n=15$ hard-regime corpus,
all $3{,}436$ canonical classes with global $L(G)=11$ admit, relative to the
recorded distinguished five-set $X$, a five-set $Y$ with $|X\cap Y|=2$ and
\[
\dd(G)-\dd(G-Y)\le 5.
\]
Among these classes, $762$ have $\dd(G)>5$ and therefore require a
nontrivial search; all $762$ pass.

The full Erd\H{o}s conjecture remains open.  The results here are
structural and computational certificates for one induction mechanism,
not a proof of the conjecture.
\end{abstract}

\section{Introduction}

For a finite simple graph $G$, define its edge-bipartization number by
\[
\dd(G)=\min\{|F|:F\subseteq E(G),\ G-F\text{ is bipartite}\}
      =|E(G)|-\mc(G).
\]
Erd\H{o}s conjectured that every triangle-free graph on $N$ vertices
satisfies
\[
\dd(G)\le \frac{N^2}{25}.
\]
The balanced blow-up of $C_5$ shows that the constant $1/25$ would be
sharp.  Balogh, Clemen and Lidick\'y proved the global bound
$N^2/23.5$ and proved the sharp $N^2/25$ bound in the two density tails
\cite{BCL}.  A recent computer-assisted preprint of Ferudun proves the
sharp value for every multiple of five up to $N=200$ \cite{Ferudun}.
Consequently, the small finite cases used below are validation and
structure-finding devices, not new extremal evaluations.

The present project asks a different question: can the sharp bound be
approached by a five-vertex induction that mirrors the extremal $C_5$
blow-up?  For a triangle-free graph on $5k$ vertices, the natural target is
as follows.

\begin{conjecture}[Candidate A]
For every triangle-free graph $G$ on $5k$ vertices, there exists a
five-set $X\subseteq V(G)$ such that
\[
\dd(G)-\dd(G-X)\le 2k-1.
\]
\end{conjecture}

If Candidate A holds for all $k$, induction from $k=1$ gives
\[
\dd(G)\le 1+3+\cdots+(2k-1)=k^2.
\]
Thus the problem becomes one of \emph{selecting} a good five-set, not only
of bounding a fixed extension.

\section{Exact finite ground truth}

All exact computations use the identity
\[
\dd(G)=|E(G)|-\mc(G).
\]
For a graph on ten vertices, global color-complement symmetry allows one
vertex to be fixed, so a direct MaxCut check uses at most $2^9$ colorings.
As a ground-truth computation, the $12{,}172$ triangle-free isomorphism
classes on ten vertices were checked exactly; the maximum is
\[
\max \dd(G)=4,
\]
attained by the balanced $C_5$ blow-up $B_2$.

This computation is useful internally because it supplies an exact base
case and a regression test for later routines.  In view of \cite{Ferudun},
it is not presented as a novelty claim.

\section{Balanced and mixed extension certificates}
\label{sec:extensions}

The first reduction studies extensions organized around the five-part
structure suggested by balanced $C_5$ blow-ups.  The implementation
separates a balanced extension model from mixed attachment patterns.
Because these are finite reduced state spaces, the conclusions in this
section are exact computational certificates for the represented models.

\subsection{Balanced-extension certificate}

For each parameter value $t=1,2,\ldots,25$, the exact C++ checker
\texttt{scripts/check\_balanced\_extension.cpp} enumerates the reduced
balanced-extension states and compares the exact deletion distance with
the target inequality.  In every batch the reported maximum gap is zero:
\[
\texttt{max\_gap}=0,
\]
and the checker reports that the inequality holds for every represented
state.

\begin{proposition}[Balanced-extension certificate]
For every balanced-extension state represented by the reduction at
$t=1,\ldots,25$, the target inequality checked by the exact program holds.
No violating state occurs in the enumerated reduced state space.
\end{proposition}

The significance of this certificate is methodological rather than a
stand-alone theorem about arbitrary graphs: it shows that the balanced
part of the proposed induction behaves exactly as the $C_5$ blow-up
heuristic predicts throughout the tested range.

\subsection{Exact mixed $t=2$ closure}

The mixed $t=2$ model is the first regime in which different attachment
types interact.  The exact closure contains
\[
4{,}013{,}521
\]
maximal-pair states, of which
\[
2{,}768{,}154
\]
are genuinely mixed.  Exact MaxCut evaluation gives maximum deletion
distance $9$ over the full state space and maximum deletion distance $8$
among the genuinely mixed states.

A separate maximal-completion monotonicity argument is used in the code:
once a legal mixed state is extended to a maximal triangle-free completion,
the completed state is the only one that must be checked for the relevant
upper-bound test.  Thus the exhaustive maximal-state computation closes
the arbitrary legal mixed $t=2$ states represented by this reduction.

\begin{proposition}[Mixed $t=2$ certificate]
Within the exact mixed $t=2$ extension model, all legal states are closed
by maximal triangle-free completion.  Among the maximal states the largest
deletion distance is $9$, while the largest genuinely mixed value is $8$.
In particular, no mixed $t=2$ obstruction to the tested induction
inequality occurs in this model.
\end{proposition}

These finite certificates motivated the next step: instead of forcing a
single extension coloring, allow the core coloring to change in a
controlled way.

\section{Coordinated flips and the selection inequality}
\label{sec:flips}

Fix a five-set $X\subseteq V(G)$ and write $H=G-X$.  Let $c$ be an optimal
2-coloring of $H$, so that the number of monochromatic edges of $H$ under
$c$ is $\dd(H)$.  Let $a$ be a 2-coloring of $X$.  Denote by $E_c(a)$ the
number of monochromatic edges contributed by $G[X]$ and the cut
$E(X,H)$ before modifying the core coloring.

For $v\in V(H)$ define
\[
w_a(v)=
\#\{\text{monochromatic edges from $v$ to $X$}\}
-
\#\{\text{bichromatic edges from $v$ to $X$}\},
\]
and
\[
\sigma_c(v)=
\#\{\text{bichromatic core edges incident with $v$}\}
-
\#\{\text{monochromatic core edges incident with $v$}\}.
\]
For $S\subseteq V(H)$ let $m_c(S)$ be the number of monochromatic edges
of $H[S]$ under $c$.  If all colors on $S$ are flipped, direct edge
accounting gives the exact gain
\[
g_c(a,S)=
\sum_{v\in S}\bigl(w_a(v)-\sigma_c(v)\bigr)
+2e(H[S])-4m_c(S).
\]
Therefore
\begin{equation}
 b_G(c^S\cup a)-\dd(H)=E_c(a)-g_c(a,S),
 \label{eq:flipidentity}
\end{equation}
where $b_G(\cdot)$ denotes the number of monochromatic edges of $G$ under
the displayed coloring.

The implemented admissible flip family consists of independent sets and
sets inducing complete bipartite subgraphs, together with their
complements.  This is deliberately narrower than the family of all
bipartite induced sets; an earlier audit that used the larger family was
corrected before the results reported here.

Let $\mu_X$ be the uniform distribution on optimal 2-colorings of $G[X]$,
counting color reversals separately.  For every fixed optimal core coloring
$c$,
\begin{equation}
\E_{a\sim\mu_X} E_c(a)
=
\dd(G[X])+\frac{e(X,H)}{2}.
\label{eq:averageextension}
\end{equation}
Define
\[
\Phi_X=
\max_{c\text{ optimal on }H}
\E_{a\sim\mu_X}
\max_{S\text{ admissible}} g_c(a,S).
\]
Combining \eqref{eq:flipidentity} and \eqref{eq:averageextension} gives the
selection bound.

\begin{proposition}[Coordinated-flip selection bound]
For every five-set $X$,
\begin{equation}
\dd(G)-\dd(G-X)
\le
\left\lfloor
\dd(G[X])+\frac{e(X,G-X)}2-\Phi_X
\right\rfloor.
\label{eq:selection}
\end{equation}
Consequently, for a graph on $5k$ vertices, Candidate A holds for $X$ if
\[
\Phi_X\ge
\dd(G[X])+\frac{e(X,G-X)}2-(2k-1).
\]
\end{proposition}

\begin{proof}
For every optimal core coloring $c$ and every assignment $a$ on $X$,
choose an admissible flip maximizing $g_c(a,S)$.  The resulting coloring
of $G$ has cost at most the right-hand side of
\eqref{eq:flipidentity}.  Since $\dd(G)$ is the minimum cost over all
2-colorings, it is at most the average of these constructed costs over
$a\sim\mu_X$.  Equation \eqref{eq:averageextension} gives the displayed
bound, and integrality gives the floor.  Finally maximize over the choice
of optimal core coloring $c$.
\end{proof}

For $k=3$, write
\[
\ell(X)=2\dd(G[X])+e(X,G-X),
\qquad
L(G)=\min_{|X|=5}\ell(X).
\]
If $L(G)\le10$, the empty flip already yields the desired bound
$\dd(G)-\dd(G-X)\le5$.  Thus the genuinely difficult $n=15$ regime begins
at $L(G)>10$.

\section{The saved $n=15$ hard regime}
\label{sec:hard}

The saved canonical corpus contains $8{,}773$ graphs with $L(G)>10$.
Their $L$-distribution is recorded by the exact audit; in particular,
exactly
\[
3{,}436
\]
classes have global $L(G)=11$.

For the $L=11$ exchange analysis, the stored distinguished five-set need
not itself minimize $\ell$.  In the $3{,}436$ records, the distinguished
sets have
\[
\ell(X)\in\{11,12,13,14,15\}
\]
with counts
\[
2502,\ 636,\ 244,\ 48,\ 6,
\]
respectively.  This distinction is important: ``global $L(G)=11$'' means
that some five-set realizes $11$, not that every later distinguished set
has $\ell=11$.

The scope of the exchange file was independently checked against the
global $L$ records.  Both contain exactly $3{,}436$ global-$L=11$ graph
IDs; the ID sets agree, there are no missing or extra records, and the
canonical graph6 strings agree record by record.

\section{Why a stronger fixed-set lemma fails}
\label{sec:failure}

The first $L=11$ audits suggested a stronger statement than Candidate A.
For a distinguished $C_5$ five-set, there are ten optimal assignments when
color reversals are counted.  Empirically, the initial saved-corpus audit
had minimum average gain $0.7$, suggesting a total gain of at least $7$
over those ten assignments.

That stronger claim is false.  A direct actual-core analysis reduces the
remaining low-imbalance configurations to $22$ signed-boundary orbits.
Exact realization and minimization show that all $22$ admit realizable
configurations with total gain below $7$.  Thus the proposed universal
``total gain at least $7$'' lemma cannot be used.

The correction is important because the stronger lemma is not the actual
selection threshold.  For a $C_5$ set $X$ in the $k=3$ case,
\eqref{eq:selection} requires
\[
\Phi_X\ge \frac{\ell(X)-10}{2}.
\]
When $\ell(X)=11$, the required average gain is only $1/2$, equivalently a
total gain of $5$ over the ten optimal assignments.  Several residual
orbits attain this smaller threshold exactly.  Therefore the $7$-gain
lemma is false while Candidate A remains intact.

The $22$ realizable residual witnesses were then checked without fixing the
original five-set.  An exact audit over all $\binom{15}{5}=3003$ deletion
sets per witness gives
\[
\max q_{\min}=4,
\qquad
q(Y)=\dd(G)-\dd(G-Y),
\]
so none is a Candidate-A obstruction.  More strikingly, restricting to
3-for-3 exchanges relative to the original five-set gives the same
minimum distribution across the $22$ witnesses:
\[
\begin{array}{c|rrrr}
\text{restricted minimum }q & 1 & 2 & 3 & 4\\ \hline
\text{number of witnesses} & 1 & 10 & 8 & 3
\end{array}
\]
Every residual witness therefore has a global minimizer obtainable by
retaining exactly two vertices of the original $X$ and replacing the other
three.  This observation motivated the full-corpus exchange test.

\section{The 3-for-3 exchange certificate}
\label{sec:3for3}

A \emph{3-for-3 exchange} of a recorded five-set $X$ is a five-set
$Y\subseteq V(G)$ satisfying
\[
|X\cap Y|=2.
\]
There are
\[
\binom52\binom{10}{3}=1200
\]
such sets when $|V(G)|=15$.

\begin{proposition}[Exact global-$L=11$ corpus certificate]
For every one of the $3{,}436$ canonical graph classes in the saved
$n=15$ hard-regime corpus with global $L(G)=11$, the recorded distinguished
five-set $X$ admits a 3-for-3 exchange $Y$ such that
\[
\dd(G)-\dd(G-Y)\le5.
\]
\end{proposition}

The complete deletion-distance distribution of the $3{,}436$ classes is
\[
\begin{array}{c|rrrrrrrr}
\dd(G) & 0&1&2&3&4&5&6&7\\ \hline
\#     & 5&49&160&472&856&1132&701&61
\end{array}
\]
For the $2{,}674$ classes with $\dd(G)\le5$, the conclusion is automatic
because $\dd(G-Y)\ge0$.  The remaining
\[
762=701+61
\]
classes are nontrivial.  The exact C++ checker searches 3-for-3 sets in a
deterministic order and passes all $762$.

For these nontrivial classes, the first successful exchange encountered by
the early-exit search has
\[
\begin{array}{c|rrr}
q(Y) & 3&4&5\\ \hline
\#   & 26&307&429
\end{array}
\]
These numbers are \emph{first-success values}, not restricted minima; no
minimality claim is attached to this distribution.

\begin{remark}
The proposition is exhaustive for the $3{,}436$ global-$L=11$ records of
the saved hard-regime corpus.  It does not, by itself, certify that the
larger saved $n=15$ corpus exhausts all triangle-free isomorphism classes
on $15$ vertices.  The scope statement is intentionally kept separate
from any stronger universe-completeness claim.
\end{remark}

The computational result suggests the following structural target.

\begin{conjecture}[3-for-3 exchange principle]
Under an appropriate structural hypothesis capturing the global-$L=11$
regime, a distinguished five-set $X$ admits a five-set $Y$ with
$|X\cap Y|=2$ and
\[
\dd(G)-\dd(G-Y)\le5.
\]
\end{conjecture}

A proof of such a statement, combined with a reduction of a minimal
Candidate-A counterexample into the relevant structural regime, would turn
the present finite certificate into a genuine induction lemma.

\section{Reproducibility and certificate scope}
\label{sec:repro}

The current computation is organized so that the numerical claims above
can be regenerated from explicit finite inputs.  The principal artifacts
are:
\begin{itemize}
\item \texttt{results/hard\_regime\_L\_n15/L\_records.tsv}, containing the
      global $L$ records;
\item \texttt{results/hard\_regime\_L\_n15/L11\_exchange.jsonl}, containing
      the $3{,}436$ exchange-analysis records;
\item \texttt{scripts/build\_L11\_3for3\_input.py}, which constructs the
      deterministic C++ input;
\item \texttt{scripts/check\_L11\_3for3\_all.cpp}, the exact 3-for-3
      checker;
\item \texttt{results/hard\_regime\_L\_n15/L11\_3for3\_input.txt} and
      \texttt{L11\_3for3\_full\_audit.tsv}, the generated input and audit
      output.
\end{itemize}

The release checker independently recomputes $\dd(G)$ using exact MaxCut,
verifies triangle-freeness, and then tests the exchange condition.  The
scope audit independently matches the global-$L=11$ IDs and graph6 strings
against the exchange file.  In the release version these inputs and outputs
should be accompanied by SHA-256 checksums and a one-command regression
script.

The distinction between three levels of evidence is maintained throughout:
(1) symbolic identities such as the flip formula,
(2) exhaustive finite certificates for explicitly defined reduced state
spaces or saved corpora, and
(3) conjectural structural statements intended for future proof.

\section{Discussion}

The computations point to a useful change in perspective.  A natural
fixed-set strengthening can fail even when the existential five-set
induction target is comfortably true.  The residual $22$-orbit analysis
shows concretely how insisting on one distinguished set can create an
artificial obstruction.  The 3-for-3 audit then shows that a small exchange
of the deletion set is sufficient across the entire saved global-$L=11$
corpus.

Two proof directions appear particularly natural.  First, an
edge-minimal counterexample to the sharp bipartization bound would have
highly constrained edge-critical behavior, potentially reducing the
number of admissible local configurations.  Second, the coordinated-flip
formula separates signed boundary imbalance from exact internal flip
corrections, providing a route to prove that one of a small family of
exchanges must have enough gain.

The immediate goal is therefore not a broader brute-force search.  It is
to convert the observed exchange phenomenon into a lemma: identify a
minimal set of structural hypotheses under which a 3-for-3 exchange must
exist, and then derive those hypotheses from a minimal counterexample to
Candidate A.

\section*{Status and non-claims}

The Erd\H{o}s triangle-free bipartization conjecture remains open in
general.  This manuscript does not claim to solve the conjecture or
JSP-000058.  The finite computations at $n=10$ and $n=15$ are not claimed
as new extremal evaluations.  The contribution is the exact certificate
framework and the structural evidence for the five-set induction and
3-for-3 exchange mechanisms.

\begin{thebibliography}{9}
\bibitem{BCL}
J. Balogh, F. C. Clemen, and B. Lidick\'y,
\emph{Max Cuts in Triangle-Free Graphs},
arXiv:2103.14179, 2021.

\bibitem{Ferudun}
A. Ferudun,
\emph{The Erd\H{o}s $n^2/25$ max-cut conjecture for small multiples of five,
via a per-root-MaxCut envelope and blow-up integrality},
arXiv:2606.28041, 2026.
\end{thebibliography}

\end{document}
